

\subsection{Monotonicity}
\label{sec:monotonicity}

In classical discovery-based querying, a result is sound the moment it is produced: as
$\Dd(T)$ grows, so does $G_{\Dd(T)}$, and BGP evaluation is monotone. Claims break this,
as discovering a composite can disregard triples that were already used. We make this
precise for the answer an engine can compute at time $T$, which can only take the claims
of retrieved composites into account.

\begin{definition}[Answer at time $T$]
\label{def:answer-at-T}
For a discovery state $W(T)$, with claims under $\alpha$, let
\[
  G^{T} \;=\; \bigcup_{d \in \Dd(T)}
  \bigl\{\, t \in \dataf(d) \;\bigm|\; d \text{ claims } t
  \ \text{ or no } c \in \Dd(T) \text{ claims } t \,\bigr\},
\]
and $\Omega(T) = \ev{B}{G^{T}}$. The answer is \emph{monotone} if
$\Omega(T) \subseteq \Omega(T')$ whenever $T \leq T'$.
\end{definition}

Once discovery is complete, $\Dd(T) = \Dd$ by Lemma~\ref{lem:frontier} and
$G^{T} = G^{\alpha}_{\Dd}$, so $\Omega(T)$ is the answer under claims. Without claims,
$G^{T} = G_{\Dd(T)}$ and $\Omega$ is monotone.

\begin{proposition}[Claims break monotonicity]
\label{prop:non-monotone}
There is a setting with times $T \leq T'$ such that $\Omega(T) \not\subseteq \Omega(T')$.
\end{proposition}

\begin{proof}
Take $B = \{(?s, p, ?o)\}$ and a composite $c$ that claims $\tau = (s, p, o)$. Let a
base source $d \notin \Dd_c$ hold $\tau$ while $\tau \notin \dataf(c)$, and let $d$ be
retrieved before $c$. If $\Dd(T) = \{d\}$, then
$\mu = \{?s \mapsto s,\ ?o \mapsto o\} \in \Omega(T)$. Once $c \in \Dd(T')$, $c$ claims
$\tau$, so $\tau$ is no longer admitted from $d$, nor from $c$ as $\tau \notin \dataf(c)$.
Hence $\mu \notin \Omega(T')$.
\end{proof}

The triple $\tau$ is claimed by $c$ but held only elsewhere; under namespace claims it is
a statement by $d$ about a resource in the namespace of $c$. It is exactly the kind of
triple the adversary of Theorem~\ref{th:undetermined} places behind
the frontier, and exactly what Theorem~\ref{th:claim-delegation} needs disregarded.
Non-monotonicity is thus the price of delegation under incomplete discovery rather than
an artefact of our definitions.~\todo{Formalise as a necessity result?} It is, however,
confined to such triples.

\begin{definition}[Covered triple]
\label{def:covered}
A triple $t$ is \emph{covered} if every composite $c \in \Dd$ that claims $t$ also holds
it, $t \in \dataf(c)$.
\end{definition}

\begin{lemma}[Covered solutions are stable]
\label{lem:covered-stable}
Under Assumption~\ref{as:disjoint}, let $\mu \in \Omega(T)$ be such that $\mu(t)$ is
covered for every $t \in B$. Then $\mu \in \Omega(T')$ for every $T' \geq T$, and
$\mu \in \ev{B}{G^{\alpha}_{\Dd}}$.
\end{lemma}

\begin{proof}
Let $t \in B$ and $\tau = \mu(t) \in G^{T}$, so some $d \in \Dd(T)$ holds $\tau$. By
Assumption~\ref{as:disjoint} at most one source claims $\tau$. If no source of
$\Dd(T')$ claims it, $\tau$ is admitted at $T'$ from $d \in \Dd(T')$. Otherwise its
claimant $c \in \Dd(T')$ holds $\tau$, as $\tau$ is covered, and $\tau$ is admitted from
$c$. Hence $\mu(B) \subseteq G^{T'}$, and $\mu \in \Omega(T')$ as
$\mathrm{dom}(\mu) = \mathrm{var}(B)$. The same argument with $\Dd$ in place of
$\Dd(T')$ gives $\mu \in \ev{B}{G^{\alpha}_{\Dd}}$.
\end{proof}

Which triples are covered depends on admissibility. For namespace claims it follows from
where a triple is stored, which Section~\ref{sec:namespace-claims} uses to resolve claim
status in practice.